% ---------------------------------------------------------------------------
% Apartat: elecció del model final, tall de períodes i reducció del model inicial
% Paquets necessaris al preàmbul:
%   \usepackage{amsmath}
%   \usepackage{booktabs}
% ---------------------------------------------------------------------------

\subsection{Elecció del model final i estratègia de períodes}

Un cop avaluat el conjunt d'inversions, s'ha adoptat com a model final la
\emph{Depuració Manual~2}, en la qual no s'ha modificat cap de les components
antidiagonals del tensor d'impedància. El criteri d'aquesta elecció no ha estat la
mera minimització de l'RMS, sinó l'equilibri global descrit anteriorment: aquest
model aconsegueix generar prou estructura per resoldre els trets geoelèctrics
d'interès i, alhora, reprodueix un resultat fortament comparable amb l'estudi de
referència, sense recórrer a l'eliminació ni a la devaluació de la informació
tridimensional continguda a les components diagonals. En altres paraules, s'ha
prioritzat un model honest i geològicament plausible per damunt d'un desajust
numèric artificialment baix.

Sobre aquesta base, s'ha decidit acotar el rang de períodes fins a $T = 1$~s. La
justificació és directa a partir de la profunditat pel·licular (\emph{skin depth}),
que en el mètode magnetotel·lúric determina la profunditat de penetració del camp
electromagnètic difós i, per tant, la profunditat màxima de sondeig efectiu:
%
\begin{equation}
    \delta \approx 503\,\sqrt{\rho\,T}\ \ [\text{m}],
    \qquad \rho\ \text{en}\ \Omega\!\cdot\!\text{m},\ T\ \text{en s},
    \label{eq:skindepth}
\end{equation}
%
on una estimació més conservadora del sondeig efectiu ve donada per la profunditat de
Niblett--Bostick, $z_{NB} \approx \delta/\sqrt{2} \approx 356\,\sqrt{\rho\,T}$. La
taula~\ref{tab:skindepth} recull tots dos valors a $T = 1$~s per a un ventall de
resistivitats representatiu del recobriment de la zona d'estudi.

\begin{table}[htbp]
    \centering
    \caption{Profunditat pel·licular $\delta$ i profunditat de sondeig de
    Niblett--Bostick $z_{NB}$ a $T = 1$~s, segons l'equació~\eqref{eq:skindepth},
    per a diferents resistivitats del recobriment.}
    \label{tab:skindepth}
    \begin{tabular}{l r r}
        \toprule
        $\rho$ ($\Omega\!\cdot\!$m) & $\delta$ a 1\,s (m) & $z_{NB}$ a 1\,s (m) \\
        \midrule
        5 \ (intrusió salina / conductor) & 1\,125 & 796 \\
        10                                & 1\,591 & 1\,125 \\
        30 \ (mescla sediment/calcària)   & 2\,755 & 1\,949 \\
        100 \ (calcària resistiva)        & 5\,030 & 3\,557 \\
        \bottomrule
    \end{tabular}
\end{table}

Com es desprèn de la taula~\ref{tab:skindepth}, fins i tot en el cas més desfavorable
per a la penetració —un recobriment fortament conductor de
$\rho \approx 5\text{–}10\ \Omega\!\cdot\!$m— a $T = 1$~s el \emph{skin depth} és de
$1{,}1\text{–}1{,}6$~km i el sondeig efectiu de Niblett--Bostick de
$\sim 0{,}8\text{–}1{,}1$~km, tots dos clarament per sota de l'objectiu de
caracterització situat a $600\text{–}700$~m, amb un factor de marge de
$\sim 1{,}6\text{–}2{,}3$. Invertint l'equació~\eqref{eq:skindepth}, per assolir
$700$~m només calen períodes de l'ordre de
$T = \rho^{-1}\,(z/503)^{2} \approx 0{,}2$~s ($\rho = 10$) o $0{,}07$~s
($\rho = 30$), de manera que l'estructura objectiu queda ben dins de la banda
mostrejada i no al límit de penetració, on quedaria mal restringida. Estendre el
rang fins a $1$~s proporciona, doncs, el coixí necessari sense excés: períodes més
llargs només afegirien sensibilitat a estructures de diversos quilòmetres,
irrellevants per a un objectiu de profunditat $\leq 700$~m.

Aquesta acotació té, a més, una conseqüència directa sobre el disseny del model. La
condició de contorn inferior del domini tridimensional, tal com estableixen
Chave i Jones, exigeix que la malla s'estengui fins a profunditats on els camps
electromagnètics siguin negligibles i s'assoleixi el règim asimptòtic, cosa que
equival a diverses profunditats pel·liculars del període més llarg al substrat
resistiu. Atès que $\delta \propto \sqrt{T}$, reduir el període màxim d'unes desenes
de segons fins a $1$~s disminueix el \emph{skin depth} màxim en el factor
corresponent (per exemple, $\sqrt{10} \approx 3{,}16$ en passar de $10$~s a $1$~s) i,
per tant, la fondària mínima que ha de tenir el model per satisfer la condició de
contorn es redueix en la mateixa proporció. Això permet aprimar el model inicial
—fent-lo compatible amb les condicions de contorn especificades per Chave i Jones—
i, alhora, reduir de manera notable el cost computacional de la inversió, sense
comprometre en cap moment la resolució de l'objectiu situat a $600\text{–}700$~m.

Convé matisar, finalment, que la profunditat pel·licular és una aproximació
unidimensional; en un medi tridimensional la sensibilitat real depèn també de la
cobertura de períodes i de la qualitat de les dades. Tanmateix, com a justificació de
primer ordre —la d'ús habitual en el disseny d'experiments magnetotel·lúrics— el
criteri és plenament vàlid i robust, atès que el marge de penetració respecte a
l'objectiu és ampli fins i tot en el cas més conductor.
